
\medskip
\phantomsection\label{proof:sga3-V-4.1-c}%
\noindent\textbf{c) Proof of \hyperref[thm:sga3-V-4.1-i]{(i)}.}
Set \(Y=\operatorname{Spec}B\) and \(p=\operatorname{Spec}i\), where
\(i:B\hookrightarrow A_0\) is the inclusion. By \hyperref[proof:sga3-V-4.1-a]{part~a)}, the morphism
\(p:X_0\to Y\) is surjective. We first show that \((Y,p)\) is a
cokernel of \((d_0,d_1)\) in the category of all ringed spaces.
Indeed, \hyperref[proof:sga3-V-4.1-b]{part~b)} shows that the set underlying \(\operatorname{Spec}B\)
is obtained from the set underlying \(X_0\) by identifying the points
\(x\) and \(y\) for which there exists \(z\in X_1\) such that
\[
  d_1z=y,\qquad d_0z=x.
\]
Moreover, since \(i\) is integral, \(p=\operatorname{Spec}i\) is
closed, and therefore \(Y\) carries the quotient topology induced by
\(X_0\).

It follows that \(p\) is open. Indeed, let \(U'\) be any open subset
of \(X_0\). Since \(d_1\) is surjective and finite locally free, it is
faithfully flat and of finite presentation, hence open. Thus
\[
  U=d_1\!\left(d_0^{-1}(U')\right),
\]
the saturation of \(U'\) for the equivalence relation defined by
\(X_\ast\), is open. Since \(Y\) has the quotient topology,
\(p(U')=p(U)\) is open.

Finally, the choice of \(B\), together with the fact that \(p,d_0\),
and \(d_1\) are affine, shows that the canonical sequence of sheaves
of rings
\[
  \mathcal O_Y\longrightarrow p_\ast\mathcal O_{X_0}
  \underset{p_\ast(\delta_0)}{\overset{p_\ast(\delta_1)}
    {\rightrightarrows}}
  p_\ast d_{0\ast}\mathcal O_{X_1}
  =p_\ast d_{1\ast}\mathcal O_{X_1}
\]
is exact.

It remains to show that \((Y,p)\) is also the cokernel of
\((d_0,d_1)\) in the category of schemes---more generally, in the
category of locally ringed spaces. Let \(q:X_0\to Z\) be a morphism of
schemes such that \(qd_0=qd_1\). By what precedes, there is a unique
morphism of ringed spaces \(r:Y\to Z\) such that \(q=rp\). We must
show that, for every \(y\in Y\), the homomorphism
\[
  \mathcal O_{Z,r(y)}\longrightarrow\mathcal O_{Y,y}
\]
induced by \(r\) is local. Since \(p\) is surjective, \(y=p(x)\) for
some \(x\in X_0\), and the assertion follows from the fact that the
homomorphism
\[
  \mathcal O_{Z,q(x)}\longrightarrow\mathcal O_{X_0,x}
\]
induced by \(q\) is local.

\medskip
\phantomsection\label{proof:sga3-V-4.1-d}%
\noindent\textbf{d) Proof of \hyperref[thm:sga3-V-4.1-ii]{(ii)}.}
This follows from \hyperref[proof:sga3-V-4.1-a]{parts~a)} and~\hyperref[proof:sga3-V-4.1-c]{c)}.

\medskip
\phantomsection\label{proof:sga3-V-4.1-e}%
\noindent\textbf{e) Proof of \hyperref[thm:sga3-V-4.1-iii]{(iii)}.}
Recall that \(\underline P\) denotes the set underlying a scheme
\(P\), and that \(\underline d:\underline P\to\underline Q\) denotes
the map induced by a morphism \(d:P\to Q\).

\medskip
\phantomsection\label{lem:sga3-V-4.1.2}%
\noindent\textbf{Lemma 4.1.2.}---%
\footnote{\label{note:sga3-V-22}Editor's note: \hyperref[lem:sga3-V-4.1.2]{this lemma}, which is used repeatedly in this
Expos\'e and in later Expos\'es (\(\mathrm{VI_A}\), \(\mathrm{VI_B}\)),
has been inserted here. It appeared as Lemma \(\mathrm{VI_B}\), 4.5.1,
in the original 1965 edition of SGA~3.}
\textit{Let \((A,\mathfrak m)\) be a local ring, let \(k\) be its
residue field, and let \(K/k\) be a field extension. There exists a
local, flat \(A\)-algebra \(B\) such that \(B/\mathfrak mB\) is
\(k\)-isomorphic to \(K\). If \(K/k\) is finite, \(B\) may moreover
be chosen finite and free over \(A\).}

This is proved in EGA \(0_{\mathrm{III}}\), 10.3.1, where it is also
shown that \(B\) may be chosen Noetherian when \(A\) is Noetherian.
For the reader's convenience, we recall the proof.

Put \(A'=A[T]\), where \(T\) is an indeterminate. If \(K=k(T)\), let
\(\mathfrak p=\mathfrak mA'\) and \(B=A'_{\mathfrak p}\). Then
\[
  B/\mathfrak mB\simeq k[T]_{(0)}=k(T).
\]
The ring \(B\) is flat over \(A'\), which is a free \(A\)-module, so
\(B\) is flat over \(A\).

Suppose next that \(K=k(t)=k[t]\), where \(t\) is algebraic over
\(k\). Set \(B=A'/(F)\), where \(F\in A'\) is a monic polynomial whose
image in \(k[T]\) is the minimal polynomial \(f\) of \(t\) over \(k\).
Then \(B\) is a free \(A\)-module of finite rank
\[
  \deg(F)=\deg(f).
\]
In particular, \(B\) is integral over \(A\), and hence every maximal
ideal of \(B\) contains \(\mathfrak m\). Since
\[
  B/\mathfrak mB\simeq k[T]/(f)\simeq K,
\]
\(B\) is local, with maximal ideal \(\mathfrak mB\). This already
proves that \(B\) can be chosen finite and free over \(A\) when
\([K:k]<\infty\).

In the general case, choose a family \((t_i)_{i\in I}\) of generators
of \(K/k\), and well-order \(I\); that is, give \(I\) a total order
\(\leq\) in which every nonempty subset has a least element. For
\(i\in I\), let \(k_i\) (respectively \(k_{<i}\)) be the subfield of
\(K\) generated by the \(t_j\) with \(j\leq i\) (respectively
\(j<i\)). By adding one element if necessary, we may suppose that
\(I\) has a greatest element \(\xi\), so that \(K=k_\xi\).

Let \(J\subseteq I\) be the set of indices \(i\) for which there
exists a direct system \((A_j)_{j\leq i}\) of local, flat
\(A\)-algebras such that
\[
  A_j/\mathfrak mA_j\simeq k_j
\]
and \(A_j\) is flat over \(A_\ell\) whenever \(\ell<j\). Suppose that
\(I\setminus J\) is nonempty, let \(i\) be its least element, and put
\[
  A'=\varinjlim_{j<i}A_j.
\]
Since tensor products commute with direct limits, \(A'\) is flat over
\(A\) and over every \(A_j\), \(j<i\), and
\[
  A'/\mathfrak mA'
  \simeq A'\otimes_A(A/\mathfrak m)
  \simeq k_{<i}.
\]
Moreover, \(A'\) is local, with maximal ideal \(\mathfrak mA'\).
Indeed, if \(x=f_j(x_j)\) is not invertible, then \(x_j\) is not
invertible and hence belongs to the maximal ideal
\(\mathfrak mA_j\) of \(A_j\); therefore \(x\in\mathfrak mA'\).

The monogenic case treated above now gives a local, flat
\(A'\)-algebra \(A_i\) such that
\[
  A_i/\mathfrak mA_i\simeq k_{<i}(t_i)=k_i.
\]
The algebra \(A_i\) is flat over every \(A_j\), \(j<i\), so \(i\in J\),
contrary to the choice of \(i\). This contradiction proves that
\(J=I\), and \(A_\xi\) has the required properties. This proves
Lemma~\hyperref[lem:sga3-V-4.1.2]{4.1.2}.
